\section{Discussion}
\label{sec:discussion}

\subsection{Key Findings and Interpretation}

\textbf{Finding 1: Alignment strategy is fingerprint-detectable from public benchmarks.}
The 83.3\% LOO-CV accuracy (p=0.031) provides evidence that DPO and SFT training produce
detectably different 4D trustworthiness profiles. Practitioners can infer a model's
likely alignment strategy from four standard benchmark evaluations available on public
leaderboards, without training-data access or model internals. The fingerprint is
sufficiently strong that 10 of 12 models are correctly identified despite subtle
differences in absolute scores (0.5--4.6pp range). For alignment auditing, this provides
an independent check when alignment provenance is incomplete or undocumented.

\textbf{Finding 2: Truthfulness, not fairness, carries the alignment signal.}
This is the paper's central surprising finding. TruthfulQA MC2 dominates discriminability
(Fisher=0.8122), while BBQ provides essentially no signal (Fisher=0.0091). DPO models
score +4.6pp higher on truthfulness (per-pair mean delta) --- the wrong direction for the fairness-driven mechanism.

We propose three competing explanations, ordered by plausibility:

\begin{enumerate}
  \item \textbf{Implicit factual reward via annotator quality ratings (most plausible).} UltraFeedback-style
    DPO training datasets explicitly rate factual quality as a preference dimension.
    If DPO models in our pool are disproportionately trained on datasets with factual quality
    ratings, preference optimization may implicitly improve truthfulness as a side effect.

  \item \textbf{Calibration improvement from contrastive training.} TruthfulQA MC2 measures
    accuracy on normalized log-probabilities. DPO's contrastive objective may improve
    calibration of the model's probability distribution independently of factual content.

  \item \textbf{Sample selection bias in community model pool.} DPO models in our sample may come
    from curators who applied additional data quality filtering independent of alignment method.
\end{enumerate}

These explanations cannot be ranked definitively from inference-only evaluation.

\textbf{Finding 3: Preference-based methods (RLHF and DPO) may share a benchmark signature.}
Llama-2-chat's misclassification as DPO suggests that the alignment fingerprint captures
a broader ``preference-based training'' cluster rather than DPO specifically. A 3-class
study (DPO vs RLHF-PPO vs SFT) is the natural extension.

\subsection{Limitations}

\textbf{L1: Causal mechanism is not established.} The empirical fingerprint is confirmed, but
the causal mechanism driving it is not. The proposed bias-avoidance mechanism is falsified;
the most plausible alternative (implicit truthfulness reward) is hypothesized but untested.
The fingerprinting claim stands independently of mechanism --- alignment strategy can be
detected without knowing why --- but the source of the fingerprint remains open.

\textbf{L2: Small sample size (n=12 models, 6 pairs).} Permutation p=0.031 is significant but
close to $\alpha$=0.05. No larger matched DPO/SFT dataset with documented alignment provenance
is publicly available without custom model training. Results are confirmatory at pilot
scale, not definitive at population scale.

\textbf{L3: 100-sample evaluation limit.} We use \texttt{--limit 100} per task for PoC efficiency.
Score estimates carry higher variance than full-task evaluation. The critical null result
($k_{BBQ}$=3/6, p=0.66) is far from threshold and robust to this; full evaluation is
immediate future work.

\textbf{L4: Community model confounds.} Not all pairs are perfectly matched for base model
and training data. Community DPO models have incompletely documented provenance.
Notably, Pair P4 uses Llama-2-7b-chat (RLHF-trained) as the SFT member in the mechanism
analysis (h-m1). Sensitivity analysis excluding P4 yields $k_{TruthfulQA}$=3/5 pairs with
DPO higher and mean delta$\approx$+3.6pp --- the qualitative conclusion (TruthfulQA dominates,
fairness null) is unchanged. Results from P4 should be interpreted cautiously as this
pair conflates RLHF vs DPO with SFT vs DPO.

\subsection{Broader Impact}

This work contributes tools for alignment auditing --- verifying how models were trained
after the fact. Positive impacts include enabling practitioners to cross-check alignment
claims against public benchmark profiles and supporting regulatory frameworks requiring
alignment documentation. Potential negative use: fingerprinting methods could be used
to identify which benchmark profile to target when misrepresenting alignment. However,
the subtlety of the signal (0.5--4.6pp differences) makes gaming this fingerprint
non-trivial. The finding that DPO fairness advantage is not supported by paired data
is potentially controversial but important for accuracy --- practitioners should not
assume DPO alignment improves fairness without empirical verification.
